\section*{Chapitre \ref{ch:ll:cpp-for-latency-i-memory} --- C++ pour la latence I~: mémoire}

\begin{solution}{exo:ll:cpp-for-latency-i-memory:1}
Décalages 0, 4, 8, 16 et 24~; les données se terminent à 26 et la taille est arrondie à 32~: 16~octets de remplissage pour 16~octets de données. Triés
(\texttt{double}, \texttt{int32}, \texttt{int16}, deux \texttt{char})~: 16~octets, sans remplissage.
\end{solution}

\begin{solution}{exo:ll:cpp-for-latency-i-memory:2}
$50\,000 \times 64 = 3\,200\,000$ octets, soit environ 3.05~MiB~: 782 pages de 4~KiB (deux \omterm{def:ll:memory-hierarchy-and-caches:tlb}{grandes pages}).
\end{solution}

\begin{solution}{exo:ll:cpp-for-latency-i-memory:3}
512~MiB représentent 131\,072 pages~; à environ \qty{1.3}{\micro\second} chacune, cela fait quelque 0.17~seconde au démarrage, sans conséquence. Sur le
chemin critique, il s'agit de \qty{1.3}{\micro\second} ajoutées au message qui touche en premier chaque page, 131\,072 fois, à des instants aléatoires.
\end{solution}

\begin{solution}{exo:ll:cpp-for-latency-i-memory:4}
\texttt{\textquotedbl{}AAPL\textquotedbl{}} tient dans le tampon interne de 15~caractères~: aucune allocation. L'identifiant de 18~caractères, lui, alloue.
\texttt{reserve(0)} n'alloue pas. \texttt{v(1)} alloue la mémoire d'un élément.
\end{solution}

\begin{solution}{exo:ll:cpp-for-latency-i-memory:5}
Remplacer le \texttt{std::function} par un appel direct ou un modèle supprime trois allocations (la copie des exécutions faite par la lambda et
la copie sur le tas, faite par la fonction, de la lambda avec ses captures)~; un vecteur de taille fixe pour les exécutions en supprime trois~; une table fixe à la place de
la table associative en supprime deux. Au total huit~; l'identifiant à taille fixe supprime le dernier.
\end{solution}

\begin{solution}{exo:ll:cpp-for-latency-i-memory:6}
Pour la queue de distribution et pour le contrôle. La médiane du tas masque ses rares chemins lents (consolidation, appels système, défauts sur les
pages nouvelles)~; une réserve n'en a aucun, ses emplacements sont contigus et pré-chargés, elle refuse plutôt que de croître, et sa capacité est un
nombre que l'équipe a choisi et qu'elle surveille.
\end{solution}

\begin{solution}{exo:ll:cpp-for-latency-i-memory:7}
Enregistrer le maximum de \texttt{live()} après chaque \texttt{create}. Le high-water mark du test de rotation est, par construction, de 10\,000
ordres vivants. Une règle de dimensionnement~: prendre le pic mesuré de l'année, le multiplier par un coefficient de sécurité pour des journées pires que toutes celles
observées (1.5 à 2), arrondir à la puissance de deux supérieure, et déclencher une alarme lorsque les objets vivants dépassent une fraction de ce nombre (disons 80~\%).
\end{solution}

\begin{solution}{exo:ll:cpp-for-latency-i-memory:8}
Les deux identifiants sont assez courts pour le tampon interne, et le même message produit à chaque fois la même clé de la table associative, de sorte qu'aucun nouveau
nœud n'est jamais créé~: le test ne peut pas voir les allocations des identifiants longs et des nouvelles clés.
Tester avec des messages de forme réaliste (identifiants longs, nouvelles clés, nombreuses exécutions) et affirmer le nombre d'allocations après la phase de chauffe.
\end{solution}

\begin{solution}{pb:ll:cpp-for-latency-i-memory:1}
\begin{enumerate}
\item $9 \times 50\,000 = 450\,000$ allocations par seconde.
\item La chaîne de l'identifiant long (1)~; le vecteur des exécutions qui croît à un, deux puis quatre éléments (3)~; le nouveau nœud de la table associative
et sa copie de la clé (2)~; la copie des exécutions faite par la lambda (1)~; le stockage sur le tas de \texttt{std::function} pour la lambda et
la copie de ses captures (2).
\item La chaîne de l'identifiant et la copie de la clé de la table~: deux.
\item La capacité du vecteur double, de sorte que le passage de quatre à huit exécutions ajoute une allocation et non quatre~; au-delà de trois, le nombre croît
seulement de façon logarithmique.
\item Environ \qty{25}{\nano\second} (une destruction et une création) sur la mesure retenue.
\item Neuf paires d'environ \qty{27}{\nano\second}~: \qty{243}{\nano\second} par message, soit quelque 12 millisecondes par seconde, environ
1~\% du thread.
\item Environ \qty{1.3}{\micro\second}~; quelque trois cents défauts donneraient \qty{400}{\micro\second}.
\item Consolider les blocs libres fragmentés~; agrandir ou réduire le tas par un appel système~; faire correspondre un grand bloc
directement, mis à zéro par le noyau~; charger les pages nouvelles en mémoire.
\item Une chaîne à taille fixe pour l'identifiant, un vecteur à taille fixe pour les exécutions, une table à entrées fixes, un calcul direct
à la place du rappel.
\item Identifiants de 24~caractères, huit exécutions, 64~entrées~; au-delà, le gestionnaire refuse (l'entrée n'est pas créée)
et la situation doit être comptée et signalée par une alarme.
\item Pour que ses pages soient chargées et sur le bon nœud, avant le début du trading~: par le thread qui les utilise, après
l'épinglage (chapitre~4).
\item 64~comparaisons de chaînes courtes, quelques dizaines de nanosecondes~; avec des milliers d'entrées, il faut une table
de hachage à adressage ouvert (chapitre~19).
\item \textbf{Résultat nommé.} Neuf allocations par message avant, aucune après~; sur ce portable, un \omterm{def:ll:cpp-for-latency-i-memory:fault}{défaut de page} coûte environ
\qty{1.3}{\micro\second}, soit quarante fois une écriture dans une page déjà présente.
\item Il remplace \texttt{operator new} par un compteur et affirme zéro allocation par message après la phase de chauffe~; une modification
qui alloue fait échouer la construction.
\item Les interruptions et l'ordonnancement (chapitre~13), les caches froids après une période calme, un verrou tenu par un autre thread, un appel
de journalisation sur le chemin critique.
\item Horodater chaque étape du message lent et corréler les messages lents avec les statistiques de l'allocateur, le compteur de défauts mineurs
(\texttt{getrusage}) et les changements de contexte de l'ordonnanceur pour le thread.
\item Refuser est une décision au coût borné et signalée par une alarme~; allouer est un coût non borné, à un moment inconnu.
\item Dans l'allocateur global de comptage du binaire de test (\texttt{firm\_arena::CountingAlloc}), et dans l'API de la réserve, qui
renvoie \texttt{None} plutôt que de croître.
\item Environ 57\,000 avec un coefficient de 1.5, arrondi à 65\,536, avec une alarme à 80~\%.
\item Décider de chaque octet que le chemin critique utilisera avant le début du trading, le toucher, et ne plus jamais demander de mémoire au tas.
\end{enumerate}
\end{solution}

\begin{solution}{iq:ll:cpp-for-latency-i-memory:1}
Le remplissage diminue, davantage d'objets tiennent dans chaque \omterm{def:ll:memory-hierarchy-and-caches:line}{ligne de cache}, et les champs lus ensemble partagent une même ligne~: moins de \omterm{def:ll:memory-hierarchy-and-caches:miss}{défauts de cache} par opération.
\iqlookfor{règles d'alignement et \omterm{def:ll:memory-hierarchy-and-caches:line}{lignes de cache}.}
\end{solution}

\begin{solution}{iq:ll:cpp-for-latency-i-memory:2}
La médiane masque les chemins lents~: consolidation, appels système pour agrandir ou réduire le tas, pages mises à zéro par le noyau, défauts
sur les pages nouvelles et verrous partagés avec d'autres threads~; leur durée ne dépend pas du programme.
\iqlookfor{queues de distribution et déterminisme plutôt que coût moyen.}
\end{solution}

\begin{solution}{iq:ll:cpp-for-latency-i-memory:3}
Un tableau d'emplacements, chacun étant une union du stockage de l'objet et d'un pointeur vers le suivant~; un pointeur de tête vers le premier emplacement libre~; create
retire un emplacement et construit l'objet sur place, destroy appelle le destructeur et le remet dans la \omterm{def:ll:cpp-for-latency-i-memory:pool}{liste libre}. Lorsque la réserve est vide, elle renvoie un pointeur nul et laisse l'appelant refuser
le travail et déclencher une alarme~; ne jamais revenir au tas sur le chemin critique.
\iqlookfor{une \omterm{def:ll:cpp-for-latency-i-memory:pool}{liste libre} intrusive et une politique d'épuisement délibérée.}
\end{solution}

\begin{solution}{iq:ll:cpp-for-latency-i-memory:4}
Les chaînes courtes sont stockées à l'intérieur de l'objet (15~caractères dans la bibliothèque GNU), les longues sur le tas~; l'allocation dépend
des données~: les tests avec des identifiants courts n'allouent jamais, la production avec des identifiants longs, si.
\iqlookfor{allocation dépendant des données.}
\end{solution}

\begin{solution}{iq:ll:cpp-for-latency-i-memory:5}
Un piège vers le noyau lorsqu'on touche une page pas encore soutenue par la mémoire~; il alloue et met à zéro une page (quelques microsecondes). Allouer
tout au démarrage, écrire dans chaque page (\omterm{def:ll:cpp-for-latency-i-memory:fault}{pré-chargement des pages}), verrouiller la mémoire contre la pagination (\texttt{mlockall}, chapitre~13) et utiliser
de \omterm{def:ll:memory-hierarchy-and-caches:tlb}{grandes pages}.
\iqlookfor{pré-chargement et verrouillage.}
\end{solution}

\begin{solution}{iq:ll:cpp-for-latency-i-memory:6}
Une construction de test qui remplace les fonctions globales d'allocation par des fonctions de comptage (ou un allocateur global de comptage en Rust),
fait tourner le composant avec des données de forme réaliste au-delà de sa phase de chauffe, et affirme que le compte ne change pas~; le test
tourne en intégration continue à chaque modification.
\iqlookfor{une assertion automatisée, avec des données réalistes.}
\end{solution}
